% ─────────────────────────────────────────────────────────────────────────────
% PixelPM — Project Management Application Report
% Compile with: pdflatex main.tex
% Required packages: graphicx, geometry, hyperref, listings, xcolor,
%                    booktabs, tabularx, float, fancyhdr, titlesec, minted (opt.)
% ─────────────────────────────────────────────────────────────────────────────
\documentclass[12pt,a4paper]{report}

% ── Page layout ──────────────────────────────────────────────────────────────
\usepackage[top=2.5cm, bottom=2.5cm, left=3cm, right=2.5cm]{geometry}

% ── Typography & encoding ────────────────────────────────────────────────────
\usepackage[T1]{fontenc}
\usepackage[utf8]{inputenc}
\usepackage{lmodern}
\usepackage{microtype}

% ── Colours ──────────────────────────────────────────────────────────────────
\usepackage{xcolor}
\definecolor{awsorange}{RGB}{255,153,0}
\definecolor{darkblue}{RGB}{35,47,62}
\definecolor{codebg}{RGB}{245,245,245}
\definecolor{codeframe}{RGB}{200,200,200}
\definecolor{linkblue}{RGB}{0,90,160}

% ── Graphics ─────────────────────────────────────────────────────────────────
\usepackage{graphicx}
\usepackage{float}
\graphicspath{{./screenshots/}{./}}  % screenshots/ for body figures; ./ for cover logo

% ── Tables ───────────────────────────────────────────────────────────────────
\usepackage{booktabs}
\usepackage{tabularx}
\usepackage{longtable}
\usepackage{array}

% ── Code listings ────────────────────────────────────────────────────────────
\usepackage{listings}
\lstset{
  backgroundcolor=\color{codebg},
  frame=single,
  framesep=4pt,
  rulecolor=\color{codeframe},
  basicstyle=\ttfamily\small,
  breaklines=true,
  breakatwhitespace=false,
  tabsize=2,
  showstringspaces=false,
  commentstyle=\color{gray},
  keywordstyle=\color{darkblue}\bfseries,
  stringstyle=\color{teal},
  numbers=left,
  numberstyle=\tiny\color{gray},
  numbersep=6pt,
  captionpos=b,
}

% ── Hyperlinks ───────────────────────────────────────────────────────────────
\usepackage[
  colorlinks=true,
  linkcolor=linkblue,
  urlcolor=linkblue,
  citecolor=linkblue,
  bookmarks=true,
  bookmarksnumbered=true
]{hyperref}

% ── Headers & footers ────────────────────────────────────────────────────────
\usepackage{fancyhdr}
\pagestyle{fancy}
\fancyhf{}
\fancyhead[L]{\small PixelPM — Project Management Application}
\fancyhead[R]{\small \leftmark}
\fancyfoot[C]{\thepage}
\renewcommand{\headrulewidth}{0.4pt}

% ── Section formatting ────────────────────────────────────────────────────────
\usepackage{titlesec}
\titleformat{\chapter}[block]
  {\normalfont\LARGE\bfseries\color{darkblue}}
  {\thechapter.}{1em}{}[\titlerule]
\titlespacing*{\chapter}{0pt}{10pt}{20pt}

% ── Misc helpers ─────────────────────────────────────────────────────────────
\usepackage{enumitem}
\usepackage{parskip}
\setlength{\parskip}{6pt}

% ─────────────────────────────────────────────────────────────────────────────
\begin{document}

% ─────────────────────────────────────────────────────────────────────────────
%  COVER PAGE  (academic style — no header/footer on this page)
% ─────────────────────────────────────────────────────────────────────────────
\begin{titlepage}
\thispagestyle{empty}
\centering

% ── Top border ───────────────────────────────────────────────────────────────
\noindent\rule{\linewidth}{2pt}\\[0.25cm]

% ── Department / School / University header ───────────────────────────────────
{\large\bfseries\MakeUppercase{Department of Computer Science and Engineering}}\\[0.3cm]
{\large\bfseries\MakeUppercase{School of Technology}}\\[0.3cm]
{\Large\bfseries\MakeUppercase{[Institution Name]}}\\[0.25cm]
{\large\bfseries SESSION 2024--25}\\[0.2cm]

\noindent\rule{\linewidth}{2pt}\\[0.6cm]

% ── Institution logo ─────────────────────────────────────────────────────────
% Institution logo — place institution_logo.png in the Report/ directory
\includegraphics[height=3.8cm]{institution_logo.png}\\[0.6cm]

\noindent\rule{\linewidth}{1pt}\\[0.5cm]

% ── Project title ─────────────────────────────────────────────────────────────
{\LARGE\bfseries PixelPM}\\[0.2cm]
{\large\textit{Project Management Application}}\\[0.15cm]
{\normalsize Technical Architecture \& Cloud Services Report}\\[0.5cm]

\noindent\rule{\linewidth}{1pt}\\[0.6cm]

% ── SUBMITTED BY ─────────────────────────────────────────────────────────────
{\large\bfseries\underline{SUBMITTED BY}}\\[0.5cm]

\begin{tabular}{@{}p{3.5cm}@{\hspace{0.3cm}}l@{\hspace{0.3cm}}l@{}}
  \textbf{NAME}        & : & Jeet Patel    \\[5pt]
  \textbf{ROLL NO.}    & : & 23BIT156      \\[5pt]
  \textbf{NAME}        & : & Manav Patel   \\[5pt]
  \textbf{ROLL NO.}    & : & 23BIT179      \\[5pt]
  \textbf{NAME}        & : & Bhumik Verma  \\[5pt]
  \textbf{ROLL NO.}    & : & 23BIT199      \\[5pt]
  \textbf{COURSE NAME} & : & PixelPM --- Project Management Application \\[5pt]
  \textbf{COURSE CODE} & : & [Course Code] \\
\end{tabular}

\vfill

% ── SUBMITTED TO ─────────────────────────────────────────────────────────────
{\large\bfseries\underline{SUBMITTED TO}}\\[0.45cm]

{\large\bfseries [Faculty Name]}\\[0.2cm]
{\normalsize [Designation]}\\[0.15cm]
{\normalsize Department of Computer Science and Engineering}\\[0.1cm]
{\normalsize [Institution Name]}\\[0.5cm]

\noindent\rule{\linewidth}{2pt}

\end{titlepage}

% ─────────────────────────────────────────────────────────────────────────────
%  TABLE OF CONTENTS
% ─────────────────────────────────────────────────────────────────────────────
\tableofcontents
\listoffigures
\clearpage

% ─────────────────────────────────────────────────────────────────────────────
\chapter{Application Overview}
% ─────────────────────────────────────────────────────────────────────────────

\section{What Is PixelPM?}

PixelPM is a full-stack, serverless \textbf{project management web application} that enables
software teams and individuals to organise, track, and collaborate on work in real time.
The application provides an intuitive Kanban-style interface alongside list and overview
perspectives, letting teams visualise progress across multiple projects simultaneously.

At its core PixelPM solves the problem of scattered task tracking by consolidating
projects, tasks, team members, and progress metrics into a single lightweight tool that
runs entirely within the AWS free tier, meaning it costs nothing to operate at small scale.

Key capabilities include:

\begin{itemize}[leftmargin=2em]
  \item \textbf{Multi-project workspace} — create, manage, and switch between independent
        projects, each with its own task board.
  \item \textbf{Task lifecycle management} — tasks move through three statuses:
        \emph{To Do}, \emph{In Progress}, and \emph{Done}, via drag-and-drop or a detail modal.
  \item \textbf{Priority tracking} — tasks carry a priority label (\emph{low}, \emph{medium},
        \emph{high}) and can be assigned to specific team members.
  \item \textbf{Three view modes} — Board (Kanban columns), List (tabular), and Overview
        (project-level summary and statistics).
  \item \textbf{AI assistant} — a floating chat widget backed by a locally-running Ollama
        \texttt{llama3.2:3b} model that is injected with a live workspace snapshot so it can
        answer questions about project status, team workloads, and task progress.
  \item \textbf{Team management} — colour-coded user avatars and per-user task assignment
        across all projects.
\end{itemize}

\section{Target Audience}

PixelPM is designed for the following user groups:

\begin{description}[leftmargin=2em, style=nextline]
  \item[Small software teams (2–15 people)]
    Teams that find enterprise tools like Jira or Asana overly complex or expensive for
    their needs. PixelPM provides the essentials—Kanban boards, priorities, assignments—
    without the configuration overhead.

  \item[Indie developers and freelancers]
    Solo practitioners managing multiple client projects simultaneously. The zero-cost
    AWS free-tier deployment means there are no ongoing hosting fees.

  \item[Bootcamp students and early-career developers]
    The open architecture and fully self-hostable local stack (Docker Compose) make it
    an ideal learning platform to study full-stack development, cloud services, and
    Infrastructure-as-Code in one coherent project.

  \item[Hackathon and side-project teams]
    Groups that need to spin up a task tracker quickly without provisioning complex
    infrastructure. A single \texttt{docker compose up} starts the full stack locally,
    and a single CloudFormation deployment creates all AWS resources in under five minutes.
\end{description}

% ─────────────────────────────────────────────────────────────────────────────
\chapter{Architecture Overview}
% ─────────────────────────────────────────────────────────────────────────────

\section{High-Level Architecture}

PixelPM follows a \textbf{serverless, three-tier architecture} split into:

\begin{enumerate}
  \item \textbf{Presentation tier} — React 19 single-page application (SPA) built with
        Vite, served via Amazon CloudFront and Amazon S3.
  \item \textbf{API tier} — Express.js RESTful backend running inside a containerised
        AWS Lambda function, exposed through Amazon API Gateway (HTTP API v2).
  \item \textbf{Data tier} — Amazon DynamoDB with three tables and one Global Secondary
        Index (GSI) for efficient project-scoped queries.
\end{enumerate}

All tiers are defined in a single \texttt{infrastructure/cloudformation.yaml}
Infrastructure-as-Code (IaC) template, making the full environment reproducible and
version-controlled.

\section{Production Request Flow}

\begin{enumerate}
  \item A user's browser loads \texttt{https://<cf-domain>} (CloudFront).
  \item CloudFront serves static React assets (\texttt{index.html}, bundled JS/CSS) from
        the private S3 bucket using Origin Access Control (OAC).
  \item API calls made by React (path prefix \texttt{/api/*}) are forwarded to the API
        Gateway origin with caching \emph{disabled}.
  \item API Gateway invokes the Lambda function via \texttt{AWS\_PROXY} integration,
        passing the full HTTP request.
  \item The Lambda (containerised Express.js) processes the request, executing DynamoDB
        operations using the AWS SDK v3.
  \item Lambda returns a response through API Gateway back to CloudFront, which forwards
        it to the browser. Execution logs are written to Amazon CloudWatch.
\end{enumerate}

\section{Local Development Flow}

For development, Docker Compose orchestrates five containers:

\begin{enumerate}
  \item \textbf{DynamoDB Local} — in-memory DynamoDB emulator on port \texttt{8000}.
  \item \textbf{db-init} — one-shot Node.js container that creates tables and seeds demo
        data (4 users, 3 projects, 8 tasks).
  \item \textbf{Ollama} — local LLM server on port \texttt{11434}.
  \item \textbf{model-init} — one-shot container that pulls \texttt{llama3.2:3b} on
        first run.
  \item \textbf{Backend} — Express.js API on port \texttt{3001}, configured with
        \texttt{DYNAMODB\_ENDPOINT} pointing to the local emulator.
  \item \textbf{Frontend} — Nginx-served React build on port \texttt{8080}, proxying
        \texttt{/api} requests to the backend container.
\end{enumerate}

% ─────────────────────────────────────────────────────────────────────────────
\chapter{System Diagram}
% ─────────────────────────────────────────────────────────────────────────────

The Mermaid diagram in \texttt{system\_diagram.mmd} (reproduced as a rendered image
below) illustrates all services, their relationships, and the two execution environments
(AWS production and local Docker).

% ── Rendered system architecture diagram ─────────────────────────────────────
% Screenshot: Rendered Mermaid system architecture diagram showing AWS services,
%             local dev containers, and all communication paths between them.
\begin{figure}[H]
  \centering
  \includegraphics[width=\textwidth]{system_diagram.png}
  \caption{PixelPM System Architecture Diagram (Production \& Local Development)}
  \label{fig:system_diagram}
\end{figure}

The diagram shows three distinct zones:

\begin{description}[leftmargin=2em, style=nextline]
  \item[Client Layer]
    The user's web browser running the React SPA.
  \item[AWS Cloud (Free Tier)]
    CloudFront CDN, S3 static hosting, API Gateway HTTP API, Lambda compute,
    three DynamoDB tables, IAM execution role, and CloudWatch logging.
  \item[Local Development Environment]
    Nginx reverse proxy, Express.js backend, DynamoDB Local emulator, and
    Ollama LLM server—all orchestrated by Docker Compose.
\end{description}

% ─────────────────────────────────────────────────────────────────────────────
\chapter{AWS Free Tier Cloud Services}
% ─────────────────────────────────────────────────────────────────────────────

PixelPM is engineered to stay within the AWS free tier at small scale. The table below
lists every AWS service used, its specific configuration in the project, and the
corresponding free-tier allowance.

\begin{longtable}{@{}p{3.2cm}p{5.4cm}p{5.6cm}@{}}
  \toprule
  \textbf{Service} & \textbf{Usage in PixelPM} & \textbf{Free-Tier Allowance} \\
  \midrule
  \endhead

  \textbf{Amazon DynamoDB} &
  Three tables: \texttt{pm\_users}, \texttt{pm\_projects}, \texttt{pm\_tasks}.
  PAY\_PER\_REQUEST billing mode. Point-in-Time Recovery enabled.
  GSI \texttt{projectId-index} on tasks table for project-scoped queries. &
  25 GB storage + 25 read/write capacity units (provisioned) or
  200 million requests/month (on-demand) — always free. \\

  \addlinespace[4pt]
  \textbf{AWS Lambda} &
  Single function \texttt{pixelpm-api}: containerised Express.js image from ECR.
  256 MB memory, 30-second timeout. Handles all API routes via
  \texttt{AWS\_PROXY} integration. &
  1,000,000 free invocations/month + 400,000 GB-seconds compute — always free. \\

  \addlinespace[4pt]
  \textbf{Amazon API Gateway (HTTP API v2)} &
  HTTP API (\texttt{pixelpm-api}) with a single catch-all route
  \texttt{ANY /\{proxy+\}} forwarding to Lambda. CORS configured for all
  origins, methods, and content-type header. Auto-deploy stage. &
  1,000,000 HTTP API calls/month free for the first 12 months. \\

  \addlinespace[4pt]
  \textbf{Amazon S3} &
  Private bucket \texttt{pixelpm-frontend-static} stores the Vite build output
  (HTML, JS, CSS, assets). All public access is blocked; CloudFront accesses
  the bucket via Origin Access Control (OAC, SigV4 signing). &
  5 GB storage + 20,000 GET + 2,000 PUT requests/month free for 12 months. \\

  \addlinespace[4pt]
  \textbf{Amazon CloudFront} &
  Single distribution with two origins: S3 (default cache behaviour — static
  assets, CachingOptimised policy) and API Gateway (\texttt{/api/*} behaviour —
  CachingDisabled policy, all HTTP methods allowed). SPA fallback returns
  \texttt{index.html} on 403/404. HTTP/2 enabled. &
  1 TB data transfer out + 10,000,000 HTTP/HTTPS requests/month — always free. \\

  \addlinespace[4pt]
  \textbf{AWS IAM} &
  Role \texttt{pixelpm-lambda-role} with managed policy
  \texttt{AWSLambdaBasicExecutionRole} (CloudWatch logs) plus an inline policy
  granting DynamoDB CRUD operations on all three tables and the tasks GSI. &
  IAM has no cost — always free. \\

  \addlinespace[4pt]
  \textbf{Amazon CloudWatch} &
  Lambda execution logs streamed automatically via the
  \texttt{AWSLambdaBasicExecutionRole} managed policy. Useful for debugging
  cold-start times, DynamoDB errors, and Ollama timeouts. &
  5 GB log data ingestion + 5 GB log storage/month free for 12 months. \\

  \bottomrule
  \caption{AWS Free-Tier Services Used in PixelPM}
  \label{tab:aws_services}
\end{longtable}

% ─────────────────────────────────────────────────────────────────────────────
\chapter{Service Integration}
% ─────────────────────────────────────────────────────────────────────────────

\section{CloudFront + S3 — Static Frontend Delivery}

The React application is compiled by Vite into a static bundle and uploaded to an S3
bucket with all public access blocked. CloudFront is the only entity that can read from
S3, enforced by an \textbf{Origin Access Control} resource that signs every S3 request
with SigV4. A bucket policy explicitly allows \texttt{s3:GetObject} only when the
requesting \texttt{AWS:SourceArn} matches the CloudFront distribution ARN.

The distribution's default cache behaviour targets the S3 origin and applies the AWS
managed \texttt{CachingOptimised} policy, providing long-lived browser caching for
hashed asset filenames. A \textbf{SPA fallback} is configured by mapping both HTTP 403
and 404 error codes to a 200 response returning \texttt{/index.html}, which is required
for React Router client-side navigation to work on direct URL access or page refresh.

\section{CloudFront + API Gateway — API Routing}

A second CloudFront cache behaviour matches the \texttt{/api/*} path pattern and routes
those requests to the API Gateway origin using the \texttt{CachingDisabled} managed
policy (no caching, all HTTP methods forwarded). This means the React frontend makes
all API calls to its own CloudFront domain with no CORS issues; CloudFront transparently
proxies them to API Gateway over HTTPS.

\section{API Gateway + Lambda — Serverless Compute}

The HTTP API uses a single \textbf{AWS\_PROXY} integration pointing to the Lambda ARN.
The route key \texttt{ANY /\{proxy+\}} is a catch-all that passes every method and path
to Lambda, which then runs the full Express.js router internally. This architecture
means the Express application requires \emph{zero code changes} to run inside Lambda —
it was designed to run as a plain Node.js process and uses the same routing logic in
both local Docker and production Lambda environments.

API Gateway auto-deploys every time the Lambda function is updated, and the stage name
\texttt{\$default} makes the base URL clean (no stage prefix in the path).

\section{Lambda + DynamoDB — Data Persistence}

The backend uses the \textbf{AWS SDK v3} DynamoDB Document Client
(\texttt{@aws-sdk/lib-dynamodb}) with \texttt{marshallOptions: \{ removeUndefinedValues: true \}}
to automatically strip undefined fields before writing to DynamoDB. The client reads
the region from \texttt{AWS\_REGION} and table names from \texttt{TABLE\_USERS},
\texttt{TABLE\_PROJECTS}, and \texttt{TABLE\_TASKS} environment variables injected by
CloudFormation at deployment time.

The tasks table has a \textbf{Global Secondary Index} named \texttt{projectId-index}
with \texttt{projectId} as the hash key and \texttt{createdAt} as the range key. This
enables \texttt{QueryCommand} to efficiently fetch all tasks belonging to a specific
project without performing a full table scan, which is critical for performance and
cost at scale:

\begin{lstlisting}[language=JavaScript, caption={GSI Query for project-scoped task retrieval (tasks.js)}]
const { Items = [] } = await docClient.send(new QueryCommand({
  TableName: TABLES.TASKS,
  IndexName: "projectId-index",
  KeyConditionExpression: "projectId = :pid",
  ExpressionAttributeValues: { ":pid": req.params.projectId },
}));
\end{lstlisting}

\section{Lambda + IAM — Least-Privilege Access}

The Lambda execution role is defined entirely within CloudFormation. It attaches the
AWS-managed \texttt{AWSLambdaBasicExecutionRole} for CloudWatch Logs access and an
inline policy granting only the six DynamoDB actions the application needs
(\texttt{GetItem}, \texttt{PutItem}, \texttt{UpdateItem}, \texttt{DeleteItem},
\texttt{Query}, \texttt{Scan}) scoped to the three table ARNs and the tasks GSI ARN.
No wildcard resource permissions are used.

\section{Lambda + CloudWatch — Observability}

All \texttt{console.log} and \texttt{console.error} calls in the Express error handler
are automatically forwarded to CloudWatch Logs via the Lambda execution environment.
This provides request-level debugging without any additional instrumentation — cold-start
durations, DynamoDB throttling errors, and upstream Ollama timeouts all surface as
log entries in the Lambda log group.

\section{AI Chat — Ollama Integration (Local Development)}

The AI route (\texttt{/api/ai}) integrates with a locally running Ollama server. Before
answering the user's question, the route handler performs three parallel DynamoDB
\texttt{ScanCommand} calls to build a live \textbf{workspace snapshot}: a structured
plain-text summary of every project, every task (with status, priority, and assignee
name), and per-user task counts. This snapshot is injected as the system prompt, so the
LLM can answer workspace-specific questions factually without hallucinating. The model
is configured with \texttt{temperature: 0.3} for more factual, less creative responses.

% ─────────────────────────────────────────────────────────────────────────────
\chapter{Application Screenshots}
% ─────────────────────────────────────────────────────────────────────────────

This chapter presents screenshots of the application user interface demonstrating
its key features and views.

\section{Dashboard — Board View (Kanban)}

% Screenshot: Main dashboard in Board/Kanban view mode showing three columns
%             (To Do, In Progress, Done) with task cards, priority badges,
%             and user avatar assignments across two or more projects.
\begin{figure}[H]
  \centering
  \includegraphics[width=\textwidth]{ui_board_view.png}
  \caption{Board view: Kanban columns with task cards, priority badges, and team avatars}
  \label{fig:ui_board_view}
\end{figure}

The Board view is the default landing screen. Each column corresponds to a task status.
Task cards display the title, assigned team member avatar, and a colour-coded priority
badge. Cards can be dragged between columns to update their status instantly via an
optimistic UI update backed by a \texttt{PUT /api/tasks/:id} call.

\section{Dashboard — List View}

% Screenshot: Dashboard in List view mode showing all tasks in a tabular/list
%             layout with columns for title, status, priority, assignee, and
%             action buttons.
\begin{figure}[H]
  \centering
  \includegraphics[width=\textwidth]{ui_list_view.png}
  \caption{List view: tabular task display with status, priority, and assignee columns}
  \label{fig:ui_list_view}
\end{figure}

The List view renders tasks in a structured table layout, making it easier to scan
many tasks at once, sort by attribute, or perform bulk reviews in sprint planning
sessions.

\section{Dashboard — Overview}

% Screenshot: Dashboard in Overview mode showing project-level summary cards
%             with task completion statistics, progress bars or percentages,
%             and team workload summaries.
\begin{figure}[H]
  \centering
  \includegraphics[width=\textwidth]{ui_overview.png}
  \caption{Overview: project-level completion statistics and team workload summary}
  \label{fig:ui_overview}
\end{figure}

The Overview presents aggregate metrics per project: total tasks, tasks done, tasks
in progress, and completion percentage. It gives project leads a quick health check
across all active projects without drilling into individual task boards.

\section{Create Project Modal}

% Screenshot: Modal dialog for creating a new project, showing input fields for
%             project name and status, with a submit/cancel button row.
\begin{figure}[H]
  \centering
  \includegraphics[width=0.7\textwidth]{ui_create_project.png}
  \caption{Create project modal: form for adding a new project to the workspace}
  \label{fig:ui_create_project}
\end{figure}

Projects are created through a modal form that POSTs to \texttt{/api/projects}. The new
project is immediately reflected in the sidebar project list and the Board view switches
to it automatically.

\section{Task Detail Modal}

% Screenshot: Task detail/edit modal showing fields for task title, status
%             selector, priority selector, assignee dropdown, and save/delete
%             action buttons.
\begin{figure}[H]
  \centering
  \includegraphics[width=0.7\textwidth]{ui_task_modal.png}
  \caption{Task detail modal: editing title, status, priority, and assignee}
  \label{fig:ui_task_modal}
\end{figure}

Clicking any task card opens a detail modal that allows editing the title, status,
priority, and assignee. Changes are persisted via \texttt{PUT /api/tasks/:id} using
a DynamoDB \texttt{UpdateCommand} that sets only the modified attributes.

\section{AI Chat Widget}

% Screenshot: Floating AI chat widget open in the bottom-right corner of the
%             screen, showing a conversation between the user and the Ollama
%             LLM assistant with a workspace-aware response about task progress.
\begin{figure}[H]
  \centering
  \includegraphics[width=0.65\textwidth]{ui_ai_chat.png}
  \caption{AI chat widget: Ollama-powered assistant with live workspace context injection}
  \label{fig:ui_ai_chat}
\end{figure}

The floating chat button in the bottom-right corner opens the AI assistant. The backend
builds a live workspace snapshot from DynamoDB and injects it into the system prompt
before forwarding the conversation to the \texttt{llama3.2:3b} model running on the
local Ollama server.

% ─────────────────────────────────────────────────────────────────────────────
\chapter{Code Overview}
% ─────────────────────────────────────────────────────────────────────────────

This chapter highlights the key source files that implement the application's core
behaviour, complemented by code screenshots.

\section{Express Server Entry Point (\texttt{backend/src/app.js})}

% Screenshot: VS Code or terminal showing the full contents of backend/src/app.js —
%             Express app setup with CORS, JSON middleware, route mounting for
%             /api/users, /api/projects, /api/tasks, /api/ai, and the error handler.
\begin{figure}[H]
  \centering
  \includegraphics[width=\textwidth]{code_backend_app.png}
  \caption{\texttt{backend/src/app.js}: Express server, route mounting, and global error handler}
  \label{fig:code_backend_app}
\end{figure}

The entry point registers CORS (origin from \texttt{CORS\_ORIGIN} environment variable),
JSON body parsing, four route modules, a health check endpoint at \texttt{/health}, and
a global error handler that normalises all thrown errors to JSON 500 responses.

\section{DynamoDB Client Configuration (\texttt{backend/src/db/dynamodb.js})}

% Screenshot: VS Code showing backend/src/db/dynamodb.js — DynamoDB client
%             construction with conditional endpoint override for local development,
%             DynamoDBDocumentClient wrapper, and exported TABLES constants.
\begin{figure}[H]
  \centering
  \includegraphics[width=\textwidth]{code_dynamodb_client.png}
  \caption{\texttt{backend/src/db/dynamodb.js}: dual-mode DynamoDB client (AWS or local emulator)}
  \label{fig:code_dynamodb_client}
\end{figure}

The client module reads \texttt{DYNAMODB\_ENDPOINT} to decide whether to point at AWS
or DynamoDB Local, making the same codebase run unchanged in both environments. Table
names are read from environment variables injected by Docker Compose or CloudFormation.

\section{Tasks Route with GSI Query (\texttt{backend/src/routes/tasks.js})}

% Screenshot: VS Code showing backend/src/routes/tasks.js — full file including
%             the ScanCommand for all tasks, QueryCommand using the projectId-index
%             GSI, PutCommand for task creation, UpdateCommand for partial updates,
%             and DeleteCommand.
\begin{figure}[H]
  \centering
  \includegraphics[width=\textwidth]{code_tasks_route.png}
  \caption{\texttt{backend/src/routes/tasks.js}: RESTful task endpoints including GSI-backed project query}
  \label{fig:code_tasks_route}
\end{figure}

The tasks router demonstrates key DynamoDB patterns: \texttt{ScanCommand} for full table
reads, \texttt{QueryCommand} against the GSI for project-scoped fetches, and a
\texttt{UpdateCommand} that dynamically builds the update expression from only the
fields present in the request body — avoiding unintentional attribute overwrites.

\section{AI Route with Workspace Context Builder (\texttt{backend/src/routes/ai.js})}

% Screenshot: VS Code showing backend/src/routes/ai.js — the buildContext() function
%             running three parallel DynamoDB scans to build a workspace snapshot,
%             the system prompt construction, and the POST /api/ai/chat handler
%             forwarding to Ollama.
\begin{figure}[H]
  \centering
  \includegraphics[width=\textwidth]{code_ai_route.png}
  \caption{\texttt{backend/src/routes/ai.js}: live workspace snapshot builder and Ollama chat proxy}
  \label{fig:code_ai_route}
\end{figure}

The AI route uses \texttt{Promise.all} to concurrently scan all three DynamoDB tables,
then assembles a structured workspace snapshot that is injected as the LLM system
prompt. This grounding technique ensures the model answers with real project data rather
than hallucinated content.

\section{React Application Component (\texttt{src/App.jsx})}

% Screenshot: VS Code showing src/App.jsx — the top-level React component with
%             useState/useEffect hooks for projects, tasks, and users state,
%             the view mode toggle logic (board/list/overview), and the drag-and-drop
%             task status update handler.
\begin{figure}[H]
  \centering
  \includegraphics[width=\textwidth]{code_app_jsx.png}
  \caption{\texttt{src/App.jsx}: centralised React state, view modes, and drag-and-drop handler}
  \label{fig:code_app_jsx}
\end{figure}

The main React component is the central hub of frontend state. It maintains projects,
tasks, users, and the active project selection using \texttt{useState} and
\texttt{useEffect} hooks. The drag-and-drop handler performs an optimistic UI update
(immediately moving the card in local state) before calling the API, so the interface
feels instant even over variable network conditions.

\section{CloudFormation Infrastructure Template (\texttt{infrastructure/cloudformation.yaml})}

% Screenshot: VS Code showing infrastructure/cloudformation.yaml — the full
%             CloudFormation template with DynamoDB table definitions, Lambda
%             function, API Gateway HTTP API, S3 bucket with OAC policy, and
%             CloudFront distribution with two origins.
\begin{figure}[H]
  \centering
  \includegraphics[width=\textwidth]{code_cloudformation.png}
  \caption{\texttt{infrastructure/cloudformation.yaml}: complete IaC definition for all AWS resources}
  \label{fig:code_cloudformation}
\end{figure}

The CloudFormation template is the single source of truth for all AWS infrastructure.
It provisions every resource — DynamoDB tables, IAM role, Lambda, API Gateway, S3,
CloudFront OAC and distribution — and outputs the CloudFront URL, API Gateway URL,
and table names for use in deployment scripts.

\section{Docker Compose Local Stack (\texttt{docker-compose.yml})}

% Screenshot: VS Code showing docker-compose.yml — all six service definitions
%             (dynamodb, db-init, ollama, model-init, backend, frontend) with
%             environment variables, depends_on conditions, and volume definitions.
\begin{figure}[H]
  \centering
  \includegraphics[width=\textwidth]{code_docker_compose.png}
  \caption{\texttt{docker-compose.yml}: six-container local development stack with health checks}
  \label{fig:code_docker_compose}
\end{figure}

The Docker Compose file orchestrates the full local stack with proper startup ordering:
DynamoDB Local must be healthy before the backend starts, the backend must pass its
health check before the frontend starts, and the model-init job ensures the
\texttt{llama3.2:3b} model is available before AI chat requests are served.

% ─────────────────────────────────────────────────────────────────────────────
\chapter{Conclusion}
% ─────────────────────────────────────────────────────────────────────────────

PixelPM demonstrates that a production-grade project management tool can be built and
operated entirely within the AWS free tier. By combining serverless compute (Lambda),
managed NoSQL storage (DynamoDB), a CDN-backed static frontend (S3 + CloudFront), and
an HTTP API (API Gateway), the application achieves high availability and low latency
with zero operational overhead and effectively zero cost at small team scales.

The architecture has several noteworthy properties:

\begin{itemize}[leftmargin=2em]
  \item \textbf{Infrastructure as Code} — every AWS resource is defined in a single
        CloudFormation template, making the environment fully reproducible and diff-able
        in version control.
  \item \textbf{Dual-mode data access} — the DynamoDB client transparently switches
        between AWS and the local emulator based on a single environment variable,
        enabling a realistic offline development experience with no code changes.
  \item \textbf{Free-tier by design} — PAY\_PER\_REQUEST DynamoDB, Lambda's generous
        free invocation quota, and CloudFront's always-free CDN tier mean a team of
        10--15 people will never exceed the free tier under normal usage.
  \item \textbf{AI without external APIs} — by routing AI chat through Ollama, the
        application preserves data privacy (workspace data never leaves the local
        machine) and incurs no per-token API costs.
  \item \textbf{Least-privilege security} — the IAM role grants exactly the six
        DynamoDB actions required, scoped to exactly the three table ARNs, with no
        wildcard permissions.
\end{itemize}

The project serves as both a practical productivity tool and a reference implementation
of modern serverless architecture patterns on AWS.

\end{document}
